\documentclass[aspectratio=169,10pt]{beamer}
% ═══════════════════════════════════════════════════════════════════════════════
% SHARED PREAMBLE — Foundation Models Course (HIT)
% To be \input by each lecture file AFTER \documentclass
% ═══════════════════════════════════════════════════════════════════════════════

% ─────────────────────────────────────────────────────────────────────────────
% BEAMER THEME & BASIC SETUP
% ─────────────────────────────────────────────────────────────────────────────
\usetheme{Madrid}
\usecolortheme{default}

% Remove navigation symbols
\beamertemplatenavigationsymbolsempty

% ─────────────────────────────────────────────────────────────────────────────
% COLORS — HIT Branding
% ─────────────────────────────────────────────────────────────────────────────
\definecolor{hitblue}{RGB}{0,51,102}        % Primary: HIT Navy
\definecolor{hitgold}{RGB}{192,148,0}       % Accent: HIT Gold
\definecolor{hitlight}{RGB}{230,240,250}    % Light background

% Override Beamer palette
\setbeamercolor{primary}{fg=white,bg=hitblue}
\setbeamercolor{secondary}{fg=hitblue,bg=hitlight}
\setbeamercolor{structure}{fg=hitblue}
\setbeamercolor{alerted text}{fg=hitgold}

% ─────────────────────────────────────────────────────────────────────────────
% PACKAGES
% ─────────────────────────────────────────────────────────────────────────────
\usepackage{amsmath}
\usepackage{amssymb}
\usepackage{mathtools}
\usepackage{booktabs}
\usepackage{graphicx}
\usepackage{hyperref}
\usepackage{tikz}
\usepackage{pgfplots}
\usepackage{tcolorbox}
\usepackage{listings}
\usepackage{xcolor}
\usepackage{algorithm}
\usepackage{algpseudocode}

% ─────────────────────────────────────────────────────────────────────────────
% TikZ LIBRARIES
% ─────────────────────────────────────────────────────────────────────────────
\usetikzlibrary{shapes,arrows,positioning,calc,chains,scopes}
\pgfplotsset{compat=1.17}

% ─────────────────────────────────────────────────────────────────────────────
% CODE LISTINGS
% ─────────────────────────────────────────────────────────────────────────────
\lstset{
  basicstyle=\ttfamily\small,
  breaklines=true,
  commentstyle=\color{gray},
  keywordstyle=\color{hitblue},
  stringstyle=\color{hitgold},
  showstringspaces=false,
  backgroundcolor=\color{hitlight},
  frame=single,
  rulecolor=\color{hitblue}
}

% ─────────────────────────────────────────────────────────────────────────────
% TCOLORBOX STYLES
% ─────────────────────────────────────────────────────────────────────────────
\tcbset{
  hitbox/.style={
    colback=hitlight,
    colframe=hitblue,
    fonttitle=\bfseries,
    boxrule=1pt
  },
  alertbox/.style={
    colback=yellow!10,
    colframe=hitgold,
    fonttitle=\bfseries,
    boxrule=1.5pt
  }
}

% ─────────────────────────────────────────────────────────────────────────────
% CUSTOM COMMANDS
% ─────────────────────────────────────────────────────────────────────────────

% Placeholder for figures
\newcommand{\placeholder}[3]{
  \begin{center}
    \fbox{\begin{minipage}[c][#2][c]{#1}
      \centering\small\color{gray}[Figure: #3]
    \end{minipage}}
  \end{center}
}

% Section divider frame
\newcommand{\sectionframe}[2]{
  \begin{frame}[plain]
    \begin{tikzpicture}[remember picture,overlay]
      \fill[hitblue] (current page.south west) rectangle (current page.north east);
      \node[anchor=center, white, font=\Large\bfseries] at (current page.center) {
        \begin{tabular}{c}
          #1 \\[0.3cm]
          {\small\color{hitgold}#2}
        \end{tabular}
      };
    \end{tikzpicture}
  \end{frame}
}

% ─────────────────────────────────────────────────────────────────────────────
% LOAD CUSTOM MACROS
% ─────────────────────────────────────────────────────────────────────────────
% ═══════════════════════════════════════════════════════════════════════════════
% SHARED MACROS — Mathematical Notation for Foundation Models Course
% ═══════════════════════════════════════════════════════════════════════════════

% ─────────────────────────────────────────────────────────────────────────────
% ACTIVATION & LAYER FUNCTIONS
% ─────────────────────────────────────────────────────────────────────────────
\newcommand{\softmax}{\mathrm{softmax}}
\newcommand{\sigmoid}{\mathrm{sigmoid}}
\newcommand{\relu}{\mathrm{ReLU}}
\newcommand{\gelu}{\mathrm{GELU}}
\newcommand{\LayerNorm}{\mathrm{LayerNorm}}
\newcommand{\FFN}{\mathrm{FFN}}
\newcommand{\MHA}{\mathrm{MHA}}
\newcommand{\Attn}{\mathrm{Attention}}

% ─────────────────────────────────────────────────────────────────────────────
% ATTENTION COMPONENTS
% ─────────────────────────────────────────────────────────────────────────────
\newcommand{\query}{Q}
\newcommand{\key}{K}
\newcommand{\val}{V}
\newcommand{\dmodel}{d_{\text{model}}}
\newcommand{\dk}{d_k}
\newcommand{\nheads}{h}

% Scaled dot-product attention formula
\newcommand{\SDPA}{
  \Attn(\query, \key, \val) = \softmax\left(\frac{\query \key^T}{\sqrt{\dk}}\right) \val
}

\newcommand{\CrossAttn}{\text{CrossAttention}}

% ─────────────────────────────────────────────────────────────────────────────
% PROBABILITY & EXPECTATION
% ─────────────────────────────────────────────────────────────────────────────
\newcommand{\E}{\mathbb{E}}
\newcommand{\Prob}{\mathbb{P}}
\newcommand{\KL}{\mathrm{KL}}
\newcommand{\ELBO}{\mathrm{ELBO}}
\newcommand{\given}{\mid}

% ─────────────────────────────────────────────────────────────────────────────
% NORMS & OPERATIONS
% ─────────────────────────────────────────────────────────────────────────────
\newcommand{\norm}[1]{\left\lVert #1 \right\rVert}
\newcommand{\abs}[1]{\left| #1 \right|}
\newcommand{\inner}[2]{\langle #1, #2 \rangle}
\newcommand{\T}{^{\top}}

% ─────────────────────────────────────────────────────────────────────────────
% VECTORS & MATRICES
% ─────────────────────────────────────────────────────────────────────────────
\newcommand{\bvec}[1]{\mathbf{#1}}
\newcommand{\mat}[1]{\mathbf{#1}}
\newcommand{\iid}{\stackrel{\text{i.i.d.}}{\sim}}
\newcommand{\defeq}{\coloneq}

% ─────────────────────────────────────────────────────────────────────────────
% LANGUAGE MODELING
% ─────────────────────────────────────────────────────────────────────────────
\newcommand{\vocab}{\mathcal{V}}
\newcommand{\context}{\text{context}}
\newcommand{\token}{t}
\newcommand{\seq}{x}
\newcommand{\ptheta}{p_\theta}
\newcommand{\pref}{p_{\text{ref}}}

% ─────────────────────────────────────────────────────────────────────────────
% RLHF & DPO
% ─────────────────────────────────────────────────────────────────────────────
\newcommand{\piref}{\pi_{\text{ref}}}
\newcommand{\pitheta}{\pi_\theta}
\newcommand{\yw}{y_w}
\newcommand{\yl}{y_l}
\newcommand{\DPOloss}{\mathcal{L}_{\mathrm{DPO}}}

% ─────────────────────────────────────────────────────────────────────────────
% RETRIEVAL & RAG
% ─────────────────────────────────────────────────────────────────────────────
\newcommand{\docs}{\mathcal{D}}
\newcommand{\qvec}{q}
\newcommand{\dvec}{d}

% ─────────────────────────────────────────────────────────────────────────────
% AGENTS & REASONING
% ─────────────────────────────────────────────────────────────────────────────
\newcommand{\obs}{o}
\newcommand{\act}{a}
\newcommand{\thought}{t}
\newcommand{\tools}{\mathcal{T}}
\newcommand{\history}{h}
\newcommand{\concat}{\oplus}

% ─────────────────────────────────────────────────────────────────────────────
% SETS & LOGIC
% ─────────────────────────────────────────────────────────────────────────────
\newcommand{\R}{\mathbb{R}}
\newcommand{\N}{\mathbb{N}}
\newcommand{\Z}{\mathbb{Z}}

\endinput


% ─────────────────────────────────────────────────────────────────────────────
% TITLE & AUTHOR (can be overridden in individual lectures)
% ─────────────────────────────────────────────────────────────────────────────
\author[D. Lederman]{Dror Lederman}
\institute{Holon Institute of Technology (HIT) \\ Data Science Department}
\date{\today}

% ─────────────────────────────────────────────────────────────────────────────
% FOOTER & HEADER
% ─────────────────────────────────────────────────────────────────────────────
\setbeamertemplate{footline}{
  \leavevmode%
  \hbox{%
    \begin{beamercolorbox}[wd=0.33\paperwidth,ht=2.25ex,dp=1ex,center]{author in head/foot}%
      \usebeamerfont{author in head/foot}\insertshortauthor
    \end{beamercolorbox}%
    \begin{beamercolorbox}[wd=0.34\paperwidth,ht=2.25ex,dp=1ex,center]{title in head/foot}%
      \usebeamerfont{title in head/foot}\insertshorttitle
    \end{beamercolorbox}%
    \begin{beamercolorbox}[wd=0.33\paperwidth,ht=2.25ex,dp=1ex,right]{frame number}%
      \usebeamerfont{frame number}\insertframenumber{} / \inserttotalframenumber\hspace*{2ex}
    \end{beamercolorbox}%
  }%
  \vskip0pt%
}

\endinput


\title{Week 10: Multi-Agent Systems}
\subtitle{Foundation Models and Agentic AI}
\author{Lecturer Name}
\institute{Holon Institute of Technology (HIT)}
\date{Semester, 2025}

\begin{document}

\begin{frame}
  \titlepage
\end{frame}

\begin{frame}{Learning Objectives}
  After this lecture you will be able to:
  \begin{itemize}
    \item Explain why \textbf{multiple agents} outperform a single agent on complex tasks.
    \item Describe \textbf{role-based}, \textbf{debate}, \textbf{delegation}, and \textbf{negotiation} patterns.
    \item Understand the \textbf{AutoGen} framework and its conversation model.
    \item Compare \textbf{CrewAI} and \textbf{LangGraph} architectures.
    \item Identify coordination challenges: communication overhead, consistency, trust.
    \item Design a multi-agent system for a given task decomposition.
  \end{itemize}
\end{frame}

\section{Why Multi-Agent?}
\sectionframe{Why Multi-Agent Systems?}{Division of labor for LLM agents}

\begin{frame}{Motivation: Why Multiple Agents?}
  \begin{columns}[T]
    \column{0.52\textwidth}
      \textbf{Limitations of single agents:}
      \begin{itemize}
        \item Context length limits task scope.
        \item A single LLM cannot excel at all sub-tasks (coding, writing, research, math).
        \item No built-in error checking or verification.
        \item Linear execution — no parallelism.
      \end{itemize}
      \vspace{0.4em}
      \textbf{Benefits of multi-agent:}
      \begin{itemize}
        \item \textbf{Specialization}: each agent optimized for its role.
        \item \textbf{Parallelism}: sub-tasks run concurrently.
        \item \textbf{Adversarial checking}: one agent critiques another.
        \item \textbf{Scalability}: add agents for larger tasks.
      \end{itemize}
    \column{0.44\textwidth}
      % tikz/multiagent_graph.tex
% Multi-agent communication graph with Orchestrator and Workers
% Usage: % tikz/multiagent_graph.tex
% Multi-agent communication graph with Orchestrator and Workers
% Usage: % tikz/multiagent_graph.tex
% Multi-agent communication graph with Orchestrator and Workers
% Usage: \input{../../tikz/multiagent_graph}
\begin{tikzpicture}[
  orch/.style={circle, draw=hitblue, fill=hitblue, text=white,
               minimum size=1.1cm, font=\small\bfseries, align=center},
  worker/.style={circle, draw=hitblue, fill=hitlightblue,
                 minimum size=0.9cm, font=\scriptsize, align=center},
  user/.style={rectangle, draw=hitgold, fill=hitgold!15, rounded corners=4pt,
               minimum width=1.4cm, minimum height=0.6cm, font=\scriptsize},
  mem/.style={cylinder, draw=gray, fill=gray!10, shape border rotate=90,
              minimum width=1.4cm, minimum height=0.7cm, font=\scriptsize},
  arr/.style={<->, >=Stealth, thick, hitblue},
  darr/.style={->, >=Stealth, thick, hitgold, dashed},
  lbl/.style={font=\scriptsize, gray},
]

% Orchestrator at center
\node[orch] (orch) {Orches-\\trator};

% Worker agents around it
\node[worker, above left=1.4cm and 1.4cm of orch] (w1) {Agent\\A};
\node[worker, above right=1.4cm and 1.4cm of orch] (w2) {Agent\\B};
\node[worker, below right=1.4cm and 1.4cm of orch] (w3) {Agent\\C};
\node[worker, below left=1.4cm and 1.4cm of orch] (w4) {Agent\\D};

% Connections: orchestrator ↔ workers
\draw[arr] (orch) -- node[above left, lbl] {task} (w1);
\draw[arr] (orch) -- node[above right, lbl] {task} (w2);
\draw[arr] (orch) -- node[below right, lbl] {result} (w3);
\draw[arr] (orch) -- node[below left, lbl] {result} (w4);

% Peer connections (debate/collaboration)
\draw[arr, dotted] (w1) -- node[above, lbl] {debate} (w2);
\draw[arr, dotted] (w3) -- node[below, lbl] {verify} (w4);

% User node
\node[user, above=1.1cm of orch] (user) {User};
\draw[darr] (user) -- node[right, lbl, pos=0.4] {goal} (orch);

% Shared memory store
\node[mem, right=2.0cm of orch] (mem) {Shared\\Memory};
\draw[darr, <->] (orch.east) -- (mem.west);
\draw[darr, bend left=15] (w2.east) to (mem.north);
\draw[darr, bend right=15] (w3.east) to (mem.south);

\end{tikzpicture}

\begin{tikzpicture}[
  orch/.style={circle, draw=hitblue, fill=hitblue, text=white,
               minimum size=1.1cm, font=\small\bfseries, align=center},
  worker/.style={circle, draw=hitblue, fill=hitlightblue,
                 minimum size=0.9cm, font=\scriptsize, align=center},
  user/.style={rectangle, draw=hitgold, fill=hitgold!15, rounded corners=4pt,
               minimum width=1.4cm, minimum height=0.6cm, font=\scriptsize},
  mem/.style={cylinder, draw=gray, fill=gray!10, shape border rotate=90,
              minimum width=1.4cm, minimum height=0.7cm, font=\scriptsize},
  arr/.style={<->, >=Stealth, thick, hitblue},
  darr/.style={->, >=Stealth, thick, hitgold, dashed},
  lbl/.style={font=\scriptsize, gray},
]

% Orchestrator at center
\node[orch] (orch) {Orches-\\trator};

% Worker agents around it
\node[worker, above left=1.4cm and 1.4cm of orch] (w1) {Agent\\A};
\node[worker, above right=1.4cm and 1.4cm of orch] (w2) {Agent\\B};
\node[worker, below right=1.4cm and 1.4cm of orch] (w3) {Agent\\C};
\node[worker, below left=1.4cm and 1.4cm of orch] (w4) {Agent\\D};

% Connections: orchestrator ↔ workers
\draw[arr] (orch) -- node[above left, lbl] {task} (w1);
\draw[arr] (orch) -- node[above right, lbl] {task} (w2);
\draw[arr] (orch) -- node[below right, lbl] {result} (w3);
\draw[arr] (orch) -- node[below left, lbl] {result} (w4);

% Peer connections (debate/collaboration)
\draw[arr, dotted] (w1) -- node[above, lbl] {debate} (w2);
\draw[arr, dotted] (w3) -- node[below, lbl] {verify} (w4);

% User node
\node[user, above=1.1cm of orch] (user) {User};
\draw[darr] (user) -- node[right, lbl, pos=0.4] {goal} (orch);

% Shared memory store
\node[mem, right=2.0cm of orch] (mem) {Shared\\Memory};
\draw[darr, <->] (orch.east) -- (mem.west);
\draw[darr, bend left=15] (w2.east) to (mem.north);
\draw[darr, bend right=15] (w3.east) to (mem.south);

\end{tikzpicture}

\begin{tikzpicture}[
  orch/.style={circle, draw=hitblue, fill=hitblue, text=white,
               minimum size=1.1cm, font=\small\bfseries, align=center},
  worker/.style={circle, draw=hitblue, fill=hitlightblue,
                 minimum size=0.9cm, font=\scriptsize, align=center},
  user/.style={rectangle, draw=hitgold, fill=hitgold!15, rounded corners=4pt,
               minimum width=1.4cm, minimum height=0.6cm, font=\scriptsize},
  mem/.style={cylinder, draw=gray, fill=gray!10, shape border rotate=90,
              minimum width=1.4cm, minimum height=0.7cm, font=\scriptsize},
  arr/.style={<->, >=Stealth, thick, hitblue},
  darr/.style={->, >=Stealth, thick, hitgold, dashed},
  lbl/.style={font=\scriptsize, gray},
]

% Orchestrator at center
\node[orch] (orch) {Orches-\\trator};

% Worker agents around it
\node[worker, above left=1.4cm and 1.4cm of orch] (w1) {Agent\\A};
\node[worker, above right=1.4cm and 1.4cm of orch] (w2) {Agent\\B};
\node[worker, below right=1.4cm and 1.4cm of orch] (w3) {Agent\\C};
\node[worker, below left=1.4cm and 1.4cm of orch] (w4) {Agent\\D};

% Connections: orchestrator ↔ workers
\draw[arr] (orch) -- node[above left, lbl] {task} (w1);
\draw[arr] (orch) -- node[above right, lbl] {task} (w2);
\draw[arr] (orch) -- node[below right, lbl] {result} (w3);
\draw[arr] (orch) -- node[below left, lbl] {result} (w4);

% Peer connections (debate/collaboration)
\draw[arr, dotted] (w1) -- node[above, lbl] {debate} (w2);
\draw[arr, dotted] (w3) -- node[below, lbl] {verify} (w4);

% User node
\node[user, above=1.1cm of orch] (user) {User};
\draw[darr] (user) -- node[right, lbl, pos=0.4] {goal} (orch);

% Shared memory store
\node[mem, right=2.0cm of orch] (mem) {Shared\\Memory};
\draw[darr, <->] (orch.east) -- (mem.west);
\draw[darr, bend left=15] (w2.east) to (mem.north);
\draw[darr, bend right=15] (w3.east) to (mem.south);

\end{tikzpicture}

  \end{columns}
\end{frame}

\section{Multi-Agent Patterns}
\sectionframe{Multi-Agent Patterns}{Orchestrator, debate, delegation, negotiation}

\begin{frame}{Pattern 1: Orchestrator-Worker}
  \begin{tcolorbox}[hitbox, title=Orchestrator-Worker Pattern]
    \textbf{Orchestrator:} receives the task, decomposes into sub-tasks, delegates to workers, synthesizes results.\\[4pt]
    \textbf{Workers:} specialized agents with a single responsibility (coding agent, research agent, QA agent).
  \end{tcolorbox}
  \vspace{0.4em}
  \textbf{Example:} Software development pipeline.
  \begin{itemize}
    \item \textit{PM agent}: writes spec.
    \item \textit{Architect agent}: designs system.
    \item \textit{Coder agent}: implements.
    \item \textit{QA agent}: writes tests.
    \item \textit{Reviewer agent}: reviews code.
  \end{itemize}
  Used by: MetaGPT~\cite{hong2023metagpt}, Devin, OpenDevin.
\end{frame}

\begin{frame}{Pattern 2: Debate}
  \begin{tcolorbox}[hitbox, title=Multi-Agent Debate (Du et al.\ 2023)]
    Two or more agents propose answers independently, then critique each other's responses.
    The debate process iterates until consensus or a judge agent decides.
    \[
      \text{Answer}_t = \text{Agent}_1(\text{Question}, \text{Answer}_{t-1}^{(2)})
    \]
  \end{tcolorbox}
  \vspace{0.4em}
  \textbf{Empirical findings:}
  \begin{itemize}
    \item Debate improves accuracy on math and factual QA (even with the same underlying model).
    \item Society-of-mind: each agent may know something others don't.
    \item Related: Constitutional AI uses a ``red team'' agent to critique responses.
  \end{itemize}
  \note{Debate is analogous to peer review or adversarial collaboration in science.}
\end{frame}

\begin{frame}{Pattern 3: Delegation and Sub-Agents}
  \begin{columns}[T]
    \column{0.55\textwidth}
      \textbf{Hierarchical delegation:}
      \begin{itemize}
        \item An agent encounters a sub-task beyond its scope.
        \item It creates or calls a \emph{sub-agent} specialized for that task.
        \item Results bubble back up the hierarchy.
      \end{itemize}
      \vspace{0.4em}
      \textbf{Dynamic vs.\ static delegation:}
      \begin{itemize}
        \item \emph{Static}: roles fixed at system design time.
        \item \emph{Dynamic}: orchestrator spawns agents on demand based on the task.
      \end{itemize}
    \column{0.41\textwidth}
      \placeholder{0.95\textwidth}{4cm}{Hierarchical agent delegation tree}
  \end{columns}
\end{frame}

\begin{frame}{Pattern 4: Negotiation and Consensus}
  \begin{tcolorbox}[hitbox]
    Agents with potentially \textbf{conflicting objectives} must reach agreement:
    \begin{itemize}
      \item Resource allocation: multiple agents request the same tool.
      \item Knowledge conflicts: agents hold contradictory beliefs.
      \item Priority conflicts: agents rank sub-tasks differently.
    \end{itemize}
    \textbf{Mechanisms:} voting, auction, turn-taking, arbitration by orchestrator.
  \end{tcolorbox}
  \vspace{0.4em}
  \textbf{Research challenge:} LLM agents can hallucinate agreed-upon decisions.
  Verification is non-trivial when all agents are language models.
\end{frame}

\section{AutoGen}
\sectionframe{AutoGen}{A framework for multi-agent conversations}

\begin{frame}{AutoGen: Architecture}
  \textbf{Wu et al.\ (2023)}~\cite{wu2023autogen} — Microsoft Research.
  \vspace{0.4em}
  \begin{tcolorbox}[hitbox, title=AutoGen Primitives]
    \begin{itemize}
      \item \textbf{ConversableAgent}: any LLM, tool-executor, or human-in-the-loop.
      \item \textbf{GroupChat}: multiple agents message each other; manager selects speaker.
      \item \textbf{UserProxyAgent}: represents the human; can execute code locally.
      \item \textbf{AssistantAgent}: LLM-powered; can generate code and plans.
    \end{itemize}
  \end{tcolorbox}
  \vspace{0.3em}
  \textbf{Key abstraction:} all agent interaction is \emph{conversational} — messages, not direct API calls.
  Human can intervene at any turn.
\end{frame}

\begin{frame}{AutoGen: Two-Agent Code Generation}
  \begin{tcolorbox}[hitbox]
    \textbf{User:} Write and run a Python script to fetch the top 10 HN stories.\\[4pt]
    \textbf{Assistant:} [generates Python code]\\[4pt]
    \textbf{UserProxy:} [executes code] \texttt{ConnectionError: rate limited}\\[4pt]
    \textbf{Assistant:} [reflects on error, generates fixed code with retry logic]\\[4pt]
    \textbf{UserProxy:} [executes fixed code] ✓ Output: [top 10 stories]
  \end{tcolorbox}
  \vspace{0.3em}
  The two-agent loop with code execution is AutoGen's most common pattern.
  The UserProxy acts as a safe execution sandbox with human approval gates.
\end{frame}

\section{LangGraph and CrewAI}
\sectionframe{LangGraph and CrewAI}{State machines and crew-based orchestration}

\begin{frame}{LangGraph: Stateful Agent Graphs}
  \begin{columns}[T]
    \column{0.55\textwidth}
      \textbf{LangGraph} (LangChain):
      \begin{itemize}
        \item Models agent workflow as a \textbf{directed graph} (nodes = LLM calls or tools).
        \item \textbf{State} is explicitly typed and passed between nodes.
        \item Supports cycles (agent loops), branching (if/else), and human-in-the-loop via interrupts.
        \item Low-level: requires explicit graph definition, but highly controllable.
      \end{itemize}
    \column{0.41\textwidth}
      \placeholder{0.95\textwidth}{3.5cm}{LangGraph: nodes (agent, tool), edges (state transitions)}
  \end{columns}
  \vspace{0.4em}
  \textbf{Use case:} When you need fine-grained control over agent control flow (e.g., customer support with escalation paths).
\end{frame}

\begin{frame}{CrewAI: Role-Based Crews}
  \begin{tcolorbox}[hitbox, title=CrewAI Concepts]
    \begin{itemize}
      \item \textbf{Agent}: has a role, goal, and backstory (embedded in system prompt).
      \item \textbf{Task}: a description + expected output + assigned agent.
      \item \textbf{Crew}: orchestrates a set of agents on a set of tasks.
      \item \textbf{Process}: sequential, hierarchical, or parallel execution.
    \end{itemize}
  \end{tcolorbox}
  \vspace{0.3em}
  \textbf{Strength:} Easy to define human-like roles (``You are a senior researcher...'').
  \textbf{Weakness:} Less flexible than LangGraph for complex conditional flows.
  \vspace{0.3em}\\
  \textbf{Framework comparison:}
  \begin{itemize}
    \item \textbf{AutoGen}: conversation-first, code-execution focus.
    \item \textbf{CrewAI}: role-first, high-level task decomposition.
    \item \textbf{LangGraph}: graph-first, fine-grained control.
  \end{itemize}
\end{frame}

\section{Coordination Challenges}
\sectionframe{Coordination Challenges}{What makes multi-agent systems hard}

\begin{frame}{Multi-Agent Failure Modes}
  \begin{tcolorbox}[alertbox, title=Common Failures]
    \begin{itemize}
      \item \textbf{Echo chambers}: agents reinforce each other's hallucinations.
      \item \textbf{Infinite loops}: agents ping-pong without progress.
      \item \textbf{Context drift}: long conversations lose early task context.
      \item \textbf{Tool conflicts}: two agents call the same tool with conflicting state.
      \item \textbf{Responsibility diffusion}: no agent takes ownership of failures.
      \item \textbf{Prompt injection propagation}: malicious input spreads across agents.
    \end{itemize}
  \end{tcolorbox}
  \vspace{0.3em}
  \textbf{Mitigations:} structured communication protocols, agent output validation,
  human checkpoints, maximum turn limits.
\end{frame}

\section{Lab / Demo}
\begin{frame}{Lab Idea: AutoGen Two-Agent Coding Pipeline}
  \begin{tcolorbox}[hitbox, title=Lab 10 — Multi-Agent Coding with AutoGen]
    \textbf{Goal:} Build a two-agent coding assistant that iteratively improves code.\\[4pt]
    \textbf{Tools:} AutoGen, Python sandbox, OpenAI API
    \begin{enumerate}
      \item Define \texttt{AssistantAgent} (LLM) and \texttt{UserProxyAgent} (code executor).
      \item Task: ``Write a Python function to scrape HN headlines, add retry logic, and save to CSV.''
      \item Observe the agent loop: code generation → execution → error → refinement.
      \item Add a third \texttt{CodeReviewerAgent} that critiques the solution before finalization.
      \item Measure: number of turns, final code quality (runs correctly?).
    \end{enumerate}
    \textbf{Deliverable:} Agent transcript + final code + turn count analysis.
  \end{tcolorbox}
\end{frame}

\section{Discussion \& Summary}
\begin{frame}{Discussion Questions}
  \begin{enumerate}
    \item Multi-agent debate improves accuracy even when the same model plays both roles. Why? Is the benefit due to more computation (multiple forward passes) or something specific to the adversarial format?
    \item AutoGen allows a human proxy to intervene at any turn. In practice, how often should a human approve agent actions? What determines the right level of autonomy?
    \item Echo chambers occur when agents reinforce each other's incorrect beliefs. How would you design a multi-agent system that is robust to this failure mode?
    \item What are the security implications of having one agent delegate tasks to sub-agents it dynamically creates? Could a compromised sub-agent affect the entire pipeline?
  \end{enumerate}
\end{frame}

\begin{frame}{Summary}
  \begin{itemize}
    \item Multi-agent systems enable \textbf{specialization, parallelism, and adversarial verification} beyond single-agent limits.
    \item Key patterns: \textbf{orchestrator-worker}, \textbf{debate}, \textbf{delegation}, \textbf{negotiation}.
    \item \textbf{AutoGen}: conversation-based, code-execution loop, human-in-the-loop checkpoints.
    \item \textbf{CrewAI}: role-based crews with high-level task assignment.
    \item \textbf{LangGraph}: explicit graph with typed state, fine-grained control flow.
    \item Coordination challenges: echo chambers, loops, drift, tool conflicts — require careful system design.
  \end{itemize}
  \vspace{0.3em}
  \textbf{Next week:} How do we measure whether any of this actually works? Evaluation of foundation models and agents.
\end{frame}

\begin{frame}[allowframebreaks]{Suggested Reading}
  \nocite{wu2023autogen, hong2023metagpt, chen2023agentverse, park2023generative}
  \bibliography{../../references}
\end{frame}

\end{document}
